\documentclass[main.tex]{subfiles}
\begin{document}

\section{Exercise 3 – Find \(V_{IN}\) for \(I_D=250\,\mu\text{A}\) (M1 active)}

\subsection*{Given}
\[
\mu_n C_{ox}=1.08\times10^{-5}\ \text{A/V}^2,\quad
\frac{W}{L}=\frac{6.29\times10^{-5}}{1\times10^{-6}}=62.9,\quad
V_T=0.6~\text{V},\quad
\lambda=0,\quad
I_D=250~\mu\text{A}.
\]

\subsection*{Method}
Using the saturation relation:
\begin{equation}
I_D=\tfrac{1}{2}\mu_n C_{ox}\tfrac{W}{L}(V_{GS}-V_T)^2,
\label{eq:3.1}
\end{equation}
the effective voltage is
\begin{equation}
V_{\text{eff}}=V_{GS}-V_T=\sqrt{\tfrac{2I_D}{\mu_n C_{ox}(W/L)}}.
\label{eq:3.2}
\end{equation}

\subsection*{Calculation}
\[
V_{\text{eff}}=\sqrt{\tfrac{2(250\times10^{-6})}{(1.08\times10^{-5})\cdot 62.9}}
=0.859~\text{V},\qquad
\boxed{V_{IN}=V_{GS}=V_T+V_{\text{eff}}=0.6+0.859=1.46~\text{V}}.
\]

\subsection*{Check}
For M1 to remain in saturation:
\[
V_{DS}\ge V_{\text{eff}}\approx0.86~\text{V}.
\]
With \(V_O\approx1.8~\text{V}\) (Exercise~4), it still remains in saturation.

\end{document}
